\section{全体输入的中心首项与统一余项}\label{gl:section}

固定方体的切向体积此前只给出测试下界。要控制任意输入，需要先把所有高水平见证
放进同一个薄层框架，再比较该薄层的质量分配。共同内核中的缺失质量可以填补，
所有见证之外的质量可以删去；这两步都增强归一化见证。剩余面层只需保留每个面的
总质量，而不要求输入的边界是图形。所得有限维缺陷关于面质量、中心和尺度联合凸，
因而可以用立方体对称群平均和 Brunn--Minkowski 不等式确定最优首项。

\begin{theorem}[全体中心方体常数的精确首项]\label{gl:main}
固定 $n\ge1$，令 $J\Subset(0,n)$ 为非空紧区间。对 $\beta\in J$ 置
\[
 c_\beta=\min(\beta,n-\beta),\qquad k_\beta=\lfloor2c_\beta\rfloor,
 \qquad L_{n,\beta}=\frac{2^n}{k_\beta!(2c_\beta)^{n-k_\beta}}.
\]
存在仅依赖 $n,J$ 的 $t_J>0$、$A_J<\infty$，使 $0<t<t_J$ 时
\begin{equation}\label{gl:uniform-asymptotic}
 \left|C_{\beta,\mathcal Q_c}(1-t)-L_{n,\beta}t^n\right|
 \le A_Jt^{n+1}\qquad(\beta\in J).
\end{equation}
特别地，任意固定 $0<\beta<n$ 都满足
\[
 \lim_{\alpha\uparrow1}\frac{C_{\beta,\mathcal Q_c}(\alpha)}{(1-\alpha)^n}
 =L_{n,\beta}.
\]
结论中的上确界涵盖所有正有限测度的 Lebesgue 可测集，允许输入随 $t$ 改变。
\end{theorem}

\subsection{共同框架与质量修改}

\begin{lemma}[全部严格见证的共同框架]\label{gl:frame}
在定理~\ref{gl:main}的参数范围内，设 $|E|=1$ 且
$H_{\beta,\mathcal Q_c}(1-t,E)\ne\varnothing$。令 $\mathscr W$ 为全部严格见证方体。
平移一个见证的中心至原点后，存在只依赖 $n,J$ 的 $B\ge1$，使充分小 $t$ 时
每个 $Q\in\mathscr W$ 均可写成
\[
 Q=tv+[-(1+ta)/2,(1+ta)/2]^n,\quad |a|+\|v\|_\infty\le B,
 \qquad Q_-\subset Q\subset Q_+,
\]
其中
$Q_\pm=[-1/2\mp Bt,1/2\pm Bt]^n$。
注意被平移的见证边长只是 $1+O_{n,J}(t)$，并未再次作伸缩。
\end{lemma}
\begin{proof}
双密度引理的尺度窗和 $\rho_{\beta/n}(1-t)=O_{n,J}(t)$ 在 $\beta\in J$ 上一致，
故所有见证的边长为 $1+O_{n,J}(t)$，且 $|Q\triangle E|=O_{n,J}(t)$。
固定一个见证 $Q_*$，则 $|Q\triangle Q_*|=O_{n,J}(t)$ 对全部 $Q\in\mathscr W$ 成立。
由引理~\ref{geo:centered-cubes}，其中心与 $Q_*$ 的中心距离为 $O_{n,J}(t)$。
用足够大的同一个 $B$ 吸收中心和边长偏差，逐坐标比较端点，即得所有包含。
这些界逐个适用于整个见证族，不需要可测选择或有限化。
\end{proof}

\begin{lemma}[填补共同核与截去外部质量]\label{gl:mass-modification}
沿用引理~\ref{gl:frame}的框架，令
$E_0=(E\cup Q_-)\cap Q_+$。则 $Q_-\subset E_0\subset Q_+$，
$|E_0|=1+O_{n,J}(t)$，且每个原严格见证仍满足
\[
 |Q|^{\beta/n-1}|E_0\cap Q|>(1-t)|E_0|^{\beta/n}.
\]
\end{lemma}
\begin{proof}
置 $\delta=|Q_-\setminus E|$、$d=|E\setminus Q_+|$。
由于 $Q_-\subset Q_*\subset Q_+$，双密度界给出 $\delta+d=O_{n,J}(t)$。
对任意原见证 $Q$，有
\[
 |E_0|=1+\delta-d>0,\qquad |E_0\cap Q|=s+\delta,\quad s=|E\cap Q|.
\]
记 $\theta=\beta/n$。函数 $u\mapsto(s+u)/(1+u)^\theta$ 在 $u\ge0$ 上递增，
因为其导数的分子是 $1+u-\theta(s+u)>0$。因此
\[
 \frac{s+\delta}{(1+\delta-d)^\theta}
 \ge\frac{s+\delta}{(1+\delta)^\theta}\ge s.
\]
乘回 $|Q|^{\theta-1}$，原见证的严格不等式便保持成立；$\delta=d=0$ 时也无须额外严格余量。
\end{proof}

\subsection{面层压缩与凸纤维}

\begin{lemma}[任意可测面层的统一上压缩]\label{gl:face-compression}
设 $Q_-\subset E_0\subset Q_+$，其中 $B$ 固定、$0<t<1/(4B)$。
存在向量 $f=(f_{j,\varepsilon})_{1\le j\le n,\,\varepsilon=\pm1}\in[-B,B]^{2n}$，
记 $e=\sum_{j,\varepsilon}f_{j,\varepsilon}$，使
\[
 |E_0|=1+te+O_{n,B}(t^2).
\]
对任意 $Q=tv+[-(1+ta)/2,(1+ta)/2]^n$ 满足 $Q_-\subset Q\subset Q_+$，
令 $g_{j,\varepsilon}=a/2+\varepsilon v_j\in[-B,B]$，则
\begin{equation}\label{gl:intersection-upper}
 |E_0\cap Q|\le1+t\sum_{j,\varepsilon}\min(f_{j,\varepsilon},g_{j,\varepsilon})
                  +O_{n,B}(t^2).
\end{equation}
余项对所有可测 $E_0$ 和全部这些 $Q$ 一致。
\end{lemma}
\begin{proof}
记 $I_t=[-1/2+Bt,1/2-Bt]$、$A_t=(1-2Bt)^{n-1}$。定义面棱柱
\[
 P_{j,\varepsilon}=\{x:\varepsilon x_j\in[1/2-Bt,1/2+Bt],\ x_k\in I_t\ (k\ne j)\}.
\]
它们和核心 $Q_-$ 的内部互不交。未覆盖的角区 $R_t$ 至少有两个坐标落在
$I_t$ 之外；$n\ge2$ 时
\[
 |R_t|\le\binom n2(4Bt)^2(1+2Bt)^{n-2}=O_{n,B}(t^2),
\]
$n=1$ 时余集为空。令
\[
 f_{j,\varepsilon}=\frac{|E_0\cap P_{j,\varepsilon}|}{tA_t}-B.
\]
因棱柱体积为 $2BtA_t$，有 $|f_{j,\varepsilon}|\le B$。
与 $Q_0$ 在核心、棱柱和角区上的质量分别比较，得到
$|E_0|=1+tA_te+O_{n,B}(t^2)=1+te+O_{n,B}(t^2)$。

在该棱柱中使用向外坐标 $u=(\varepsilon x_j-1/2)/t$。
由于 $Q$ 包含核心，其余坐标的整个横截面都在 $Q$ 中；法向允许区间恰为
$-B\le u\le g_{j,\varepsilon}$。于是
\[
 |E_0\cap Q\cap P_{j,\varepsilon}|
 \le\min\{|E_0\cap P_{j,\varepsilon}|,|Q\cap P_{j,\varepsilon}|\}
 =tA_t[B+\min(f_{j,\varepsilon},g_{j,\varepsilon})].
\]
加上共同核心和角区的统一体积界，再与 $|Q_0|=1$ 比较，即得
式~\eqref{gl:intersection-upper}。整个论证只用了质量和容积，无边界正则性假设。
\end{proof}

定义联合凸缺陷
\begin{align}
 D_f(v,a)&=(n-\beta)a+\theta e-
          \sum_{j,\varepsilon}\min(f_{j,\varepsilon},a/2+\varepsilon v_j)\notag\\
 &=(n-\beta)a-(1-\theta)e+
          \sum_{j,\varepsilon}(f_{j,\varepsilon}-a/2-\varepsilon v_j)_+,
 \qquad\theta=\beta/n,\label{gl:defect}
\end{align}
并令 $K_f(\lambda)=\{v:\exists a\in\mathbb R,\ D_f(v,a)<\lambda\}$。

所需标准工具是 Brunn--Minkowski 的凸体形式；见 Gardner~\cite{Gardner2002}。
\begin{classical}[Brunn--Minkowski 不等式]\label{gl:bm}
若 $K,L\subset\mathbb R^n$ 为非空有界开凸集，$0<s<1$，则
\[
 |(1-s)K+sL|^{1/n}\ge(1-s)|K|^{1/n}+s|L|^{1/n}.
\]
这由凸体版本对闭包应用得到；有限多个集合的加权版本由归纳得到。
\end{classical}

\begin{lemma}[面质量纤维的最优体积]\label{gl:fibre-volume}
对任意 $f\in\mathbb R^{2n}$、$\lambda>0$，$K_f(\lambda)$ 为有界开凸集或空集，且
\begin{equation}\label{gl:fibre-bound}
 |K_f(\lambda)|\le L_{n,\beta}\lambda^n.
\end{equation}
\end{lemma}
\begin{proof}
$D_f$ 关于 $(f,v,a)$ 联合凸，故投影 $K_f(\lambda)$ 凸；它是固定 $a$ 时关于 $v$ 的开集的并，
因而开。并且
\[
 D_f\ge(n-\beta)a-(1-\theta)e,\qquad D_f\ge-\beta a+\theta e,
 \qquad |D_f-D_0|\le2\|f\|_1.
\]
前两式保证关于 $a$ 的强制性。对 $a\ge0$，
$D_0=(n-\beta)a+\sum_j(|v_j|-a/2)_+$；对 $a=-u\le0$，
$D_0=\beta u+\sum_j(|v_j|-u/2)_+$。由引理~\ref{tan:parameter-elimination}，
\[
 \inf_aD_0(v,a)=N_{c_\beta}(v).
\]
因而 $K_f(\lambda)\subset\{N_{c_\beta}<\lambda+2\|f\|_1\}$，这给出有界性。

若纤维为空，结论显然。否则令 $G$ 为坐标置换和符号翻转组成的有限群。
它同时置换面参数并正交作用于中心，满足
$D_{gf}(gv,a)=D_f(v,a)$，故 $K_{gf}=gK_f$ 且各纤维体积相等。
联合凸性给出
\[
 \frac1{|G|}\sum_{g\in G}K_{gf}(\lambda)\subset K_{\bar f}(\lambda),
 \qquad \bar f_{j,\varepsilon}=e/(2n).
\]
对每个参与平均的中心分别取一个严格的尺度见证，即可保证包含中的严格不等式。
由定理~\ref{gl:bm}，左侧体积不小于 $|K_f(\lambda)|$。
最后，直接代入式~\eqref{gl:defect}得到
\[
 D_{\bar f}(v,a)=D_0(v,a-e/n).
\]
所以 $K_{\bar f}(\lambda)=K_0(\lambda)=\lambda S_{n,c_\beta}$，体积恰为
$L_{n,\beta}\lambda^n$。这证明式~\eqref{gl:fibre-bound}。
\end{proof}

\begin{proposition}[半阶面模型的非刚性与严格阈值]\label{gl:flat-fibre}
当 $n=2$、$\beta=1$，取
$f_{1,+}=f_{1,-}=\delta$、$f_{2,+}=f_{2,-}=-\delta$，$\delta\ge0$。
则
\[
 \inf_aD_f(v,a)=\max\{2\delta,|v_1|+|v_2|\}.
\]
因此 $0\le\delta<1/2$ 时 $K_f(1)=\{|v_1|+|v_2|<1\}$、体积为 $2$；
$\delta\ge1/2$ 时严格纤维为空。
\end{proposition}
\begin{proof}
令 $b=a/2$，由 $x_+=(x+|x|)/2$ 及相反面的配对，有
\[
 D_f(v,a)=\max(|\delta-b|,|v_1|)+\max(|-\delta-b|,|v_2|).
\]
写 $p=|v_1|$、$q=|v_2|$。该式总不小于 $2\delta$ 和 $p+q$。
若 $p+q\ge2\delta$，区间 $[\delta-p,\delta+p]$ 与
$[-\delta-q,-\delta+q]$ 相交，在交集中选 $b$ 即得值 $p+q$。
若 $p+q\le2\delta$，取 $b\in[-\delta+q,\delta-p]$，两距离之和等于
$2\delta$ 且分别不小于 $p,q$，故达到值 $2\delta$。这证明所述下确界及严格水平集结论。
\end{proof}

命题~\ref{gl:flat-fibre}还说明，面参数收敛并不保证严格纤维体积连续。
因此下一步直接比较带余量的 $K_f(1+A_{n,J}t)$，不把它替换为极限参数处的 $K_f(1)$。

\subsection{统一上界的闭合}
\begin{proof}[定理~\ref{gl:main}的证明]
尺度不变性允许令 $|E|=1$。光晕为空时上界成立；否则应用共同框架及质量修改引理。
每个原严格见证仍以 $|E_0|$ 归一化严格成立。在面质量参数 $f$ 下，
由式~\eqref{gl:intersection-upper}，一致地有
\begin{align*}
 (1+ta)^{\beta-n}\frac{|E_0\cap Q|}{|E_0|^\theta}
 &\le(1+ta)^{\beta-n}
 [1+t\sum_{j,\varepsilon}\min(f_{j,\varepsilon},g_{j,\varepsilon})+O_{n,J}(t^2)]
 [1+te+O_{n,J}(t^2)]^{-\theta}\\
 &=1-tD_f(v,a)+O_{n,J}(t^2).
\end{align*}
这里所有参数都在同一有界区间，质量分母保持正的统一下界，所以 Taylor 余项
与输入及见证无关。结合严格下界 $1-t$，得到 $D_f(v,a)<1+A_{n,J}t$。
中心方体的见证中心就是光晕取值点，所以逐点存在这样的尺度参数，整个原光晕满足
\[
 t^{-1}H_{\beta,\mathcal Q_c}(1-t,E)\subset K_f(1+A_{n,J}t).
\]
引理~\ref{gl:fibre-volume}给出
\[
 |H_{\beta,\mathcal Q_c}(1-t,E)|\le
 L_{n,\beta}(1+A_{n,J}t)^n t^n.
\]
对所有输入取上确界，即得定理的统一上界。

下界取 $E=Q_0$。定理~\ref{hd:all-orders}和命题~\ref{hd:explicit-remainder}
的 $N=1$ 情形在同一个 $J$ 上给出
$|H_{\beta,\mathcal Q_c}(1-t,Q_0)|=L_{n,\beta}t^n+O_{n,J}(t^{n+1})$。
其 $N=1$ 积分主项与 $t$ 无关；由定理~\ref{tan:unit-cube-limit}逐参数识别为 $L_{n,\beta}$。
此处的误差使用跨参数一致截断，而非仅逐参数的次首项公式。
最后 $L_{n,\beta}$ 在 $J$ 上有界，故两侧合并得到式~\eqref{gl:uniform-asymptotic}。
\end{proof}

\begin{corollary}[全体余项的sharp阶：正次首项区]\label{gl:positive-remainder}
固定 $\min(1,n/2)\le\beta<n$。则
\[
 C_{\beta,\mathcal Q_c}(1-t)-L_{n,\beta}t^n\asymp_{n,\beta}t^{n+1}
 \qquad(t\downarrow0).
\]
\end{corollary}
\begin{proof}
定理~\ref{hd:second-coefficient}在此范围给出 $T>0$，故固定方体下界使左侧
至少为 $(T/2)t^{n+1}$。定理~\ref{gl:main}给出同阶上界。
\end{proof}

\begin{remark}[首项之后的问题]\label{gl:next-order}
本定理确定全体集合的最优首项，但未确定 $t^{n+1}$ 的全体最优系数。
面层压缩留下的角区误差恰会影响这一阶，因而固定方体的次首项 $T$ 不能直接传递。
即使有限维纤维达到首项体积，也不必迫使所有面参数相同；后续需要针对这些并列
近极值构型研究二阶损失，而不预设完整高阶幂级数存在。
\end{remark}
